% Proven local refinement; proposed full Selmer refactor explicitly separated.
\subsection{The mirror injection has a symmetric image}
\label{sec:symmetric-mirror}

The source actually proves more than its stated global-length bound uses.
For a finite module $W=T/\pi^nT$, let
\[
 \mathcal C_n^+(T)=
 \{\eta\in\operatorname{Hom}_{G_K}(W,W^{*c}):
          \eta(u)(v)=\eta(v)(u)\text{ for all }u,v\in W\}.
\]
The superscript $+$ refers to symmetry of the bilinear form, not to the
action of complex conjugation on the Selmer group.

\begin{proposition}[Sign-sensitive mirror bound]
\label{prop:symmetric-mirror}
Under the hypotheses used to construct the source's mirror injection,
its image lies in $\mathcal C_n^+(T)$.  Consequently
\[
 \operatorname{length}_{\mathcal O}\mathcal S_n^{\mathrm{dual}}
 \le\operatorname{length}_{\mathcal O}\mathcal C_n^+(T).
\]
This strengthening uses the already-proved identity
\texttt{mirrorForm\_eq\_zero}, together with the spanning hypothesis;
it does not require a division by $2$.
\end{proposition}
\begin{proof}
Write $H$ for the subgroup on which the source detects cocycles.
Its construction has $z(H)$ spanning $W$, and
\[
 \Phi_n(y)(z(h))=y(h^c)\quad(h\in H).
\]
The identity \texttt{mirrorForm\_eq\_zero} states
\[
 y(h')(t z(h))-y(h^c)(t z(h'^c))=0
 \qquad(h,h'\in H,\ t\in\mathcal O).
\]
Put $h'=k^c$.  Since $H$ is stable under conjugation, this gives
\[
 \Phi_n(y)(z(k))(t z(h))
 =\Phi_n(y)(z(h))(t z(k)).
\]
Extend in both arguments using the spanning of $z(H)$ and
$\mathcal O$-bilinearity.  This proves symmetry on all of $W\times W$.
The source already proves $\Phi_n$ injective, so the length inequality
follows.
\end{proof}

Define the smaller global term
\[
 C^+(T)=\sup_{n\ge1}\operatorname{length}_{\mathcal O}\mathcal C_n^+(T).
\]
Appending the symmetry equations $X-X^{\mathsf t}=0$ to the finite
intertwining matrix gives the following exact criterion:
\[
 C^+(T)<\infty
 \quad\Longleftrightarrow\quad
 \text{there is no nonzero symmetric rational conjugate Tate-dual map.}
\]
Indeed, the Smith formula proved above applies to the augmented relation
matrix.  Anticyclotomic twists preserve these symmetry equations as well
as the intertwining equations, so they preserve this refined obstruction.

\begin{proposition}[The opposite sign admits a bounded mirror target]
\label{prop:opposite-sign}
Suppose the rational conjugate Tate-dual Hom space is one-dimensional,
generated by an alternating pairing.  Let $P$ be a primitive integral
generator of its lattice, and let $a_1,\dots,a_r$ be the nonzero Smith
valuations of the full intertwining matrix.  Put
$\gamma_n=\sum_j\min(n,a_j)$ and
$d_2=v_\pi(2)$.  Then
\[
 \operatorname{length}\mathcal C_n^+(T)
 \le \gamma_n+\min(n,d_2)
 \le\sum_j a_j+d_2.
\]
For odd residue characteristic the term involving $2$ vanishes.
\end{proposition}
\begin{proof}
The quotient of the full finite-level Hom module by
$(\mathcal O/\pi^n)P$ has length $\gamma_n$.
The kernel of its restriction to the symmetric submodule consists of
the scalar multiples $bP$ that are symmetric modulo $\pi^n$.
Since $P^{\mathsf t}=-P$ and $P$ is primitive, this is equivalent to
$2b=0$ in $\mathcal O/\pi^n$.
That annihilator has length $\min(n,d_2)$.  Apply the length formula
to the kernel and image of the restricted quotient map.
\end{proof}

The source has an additional integral irreducibility hypothesis:
$\operatorname{span}_{\mathcal O}\rho(G)=M_d(\mathcal O)$.
Under it, the sign dichotomy becomes exact.

\begin{proposition}[Exact dichotomy under the source's matrix-span hypothesis]
\label{prop:exact-sign}
Assume $\operatorname{span}_{\mathcal O}\rho(G)=M_d(\mathcal O)$ and
$\operatorname{Hom}_{F[G]}(V,D(V))\ne0$.  A primitive integral generator
$P$ of this Hom space is invertible over $\mathcal O$, and has a sign
$P^{\mathsf t}=\epsilon P$, $\epsilon\in\{1,-1\}$.
For every $n$ the full finite-level Hom module is exactly
$(\mathcal O/\pi^n)P$.  Consequently
\[
 \operatorname{length}\mathcal C_n^+(T)=
 \begin{cases}
  n,&\epsilon=1,\\
  \min(n,v_\pi(2)),&\epsilon=-1.
 \end{cases}
\]
In the second case $C^+(T)=v_\pi(2)$, which is zero at odd $\ell$;
its normalization by $v_\ell(\ell)=1$ is at most $1$ at $\ell=2$.
\end{proposition}
\begin{proof}
Reducing the matrix-span condition modulo $\pi$ shows that the residual
representation is absolutely irreducible.  A primitive integral
intertwiner $P$ reduces to a nonzero map between the residual
representation and its conjugate Tate dual.  Its kernel is invariant
under every matrix, and hence is zero.  Therefore its reduction is
invertible and $P\in\mathrm{GL}_d(\mathcal O)$.
After identifying the target by $P$, an equivariant map modulo $\pi^n$
is a matrix commuting with $\rho(G)$, hence with all of
$M_d(\mathcal O/\pi^n)$.  Such a matrix is scalar, as is seen by commuting
with all matrix units.  This proves the asserted form of the Hom module.
Transposition preserves the conjugate-duality equivariance equations:
the equivalent bilinear equation is
$B(\rho(g)u,\rho(g^c)v)=\chi(g)B(u,v)$, with $\chi(g^c)=\chi(g)$.
Thus $P^{\mathsf t}=\epsilon P$ for a scalar, and transposing twice gives
$\epsilon^2=1$.  In characteristic zero its sign is $1$ or $-1$.
Finally $bP$ is symmetric if and only if $(1-\epsilon)b=0$ modulo
$\pi^n$.  This gives the two lengths.
\end{proof}

For the Tate module of an abelian variety descending to a totally real
field, the relevant pairing has the other sign.  If $e$ is the
alternating Weil pairing and $C$ is complex conjugation, then
$e(Cv,Cw)=-e(v,w)$.  Therefore
\[
 B(w,v)=e(w,Cv)=-e(Cw,v)=e(v,Cw)=B(v,w).
\]
The forced polarization map is symmetric.  It survives precisely the
symmetry condition of Proposition~\ref{prop:symmetric-mirror}, so neither
odd $\ell$ nor removal of $2$-torsion eliminates the BSD obstruction.

\paragraph{Why symmetry alone cannot justify quotienting.}
Let $R=\mathcal O/\pi^n$, $W=R$, and $H=(R,+)$, with conjugation the
identity.  Put $z(h)=h$ and $y_b(h)(v)=bhv$ (composed with a
generating additive character if one uses character-module values).  Then $z(H)$
spans $W$ and the mirror expression is identically zero.  The associated
map is $\Phi(b)=b\operatorname{id}$, which is injective; its image is
exactly the genuine symmetric pairing line.  Quotienting by that line
kills the whole image of length $n$.  This is a linear-algebra
counterexample to a proposed deduction from mirror symmetry and spanning;
it is not asserted to be a Galois Selmer datum.

\paragraph{A precise candidate refactor.}
Replace the field \texttt{congruenceTerm\_ne\_top} of
\texttt{SelmerDatum} by finiteness of $C^+(T)$ and define
$\beta^+(T)=C^+(T)+\text{local terms}$.
The strengthened mirror inequality then supplies the required endpoint
of the descent.  In the inspected source, subsequent uses are:
\begin{enumerate}
\item \texttt{Descent.lean}: the dual-length estimate is inserted into the
      tied Selmer-length estimate;
\item \texttt{Line.lean}: these lengths supply the constant-term
      polynomial bound and the finite budget for \texttt{IsTorsionLevel};
\item \texttt{GeneralTheorem.lean}: the finite budget is used for
      torsion-point specialization and the simultaneous twist theorem.
\end{enumerate}
The old source data automatically satisfy the proposed new condition
because $C^+(T)\le C(T)$.  The prospective enlargement includes
opposite-sign self-dual representations when all the other scalar,
irreducibility, local-rank, and ramification assumptions are established.

\emph{Status.} The symmetric-image theorem, the Smith criterion, the
opposite-sign bound, and the polarized counter-obstruction are proved
above.  The full generalized twisting theorem is a proposed source
refactor, supported by the displayed dependency chain and usage audit;
it has not been formally rebuilt here.  No instance of a new arithmetic
vanishing theorem is claimed.  For symmetric polarized data the
uncontrolled term $\Phi_n^{-1}(H/\pi^nH)$ remains an explicit arithmetic
obligation.
